\section{Experimental design}
\label{sec:design}

Studies~1 and~2 follow protocols frozen before their runs and tagged in the repository (\texttt{protocol-v1},
amended as \texttt{protocol-v2}; Study~2 as \texttt{protocol-v3}). Every change made after freezing is listed,
dated and explained in the protocol's ``Deviations'' section and summarised in Appendix~\ref{app:deviations}.

\paragraph{Cases.} Five agents, all with \texttt{gpt-5.6-terra} at default parameters, in a loop of our own
(OpenAI Responses API with function calling), because hosted tool integrations do not expose tool outputs to the
client.
\begin{itemize}
\item \textbf{C1, editorial plan}: twelve tools of a search-data API (keyword research, trends, search results,
  rankings, page parsing) reached through its MCP server; the agent plans articles for the author's own website.
\item \textbf{C2, content gaps}: search results and page parsing; the agent compares what ranks on Google with
  what the website has published.
\item \textbf{C3, literature scan}: two tools on the OpenAlex catalogue (CC0), five results per search.
\item \textbf{C4, billing support}$^\dagger$: seven tools of a simulated MCP server whose payloads follow the
  public OpenAPI description of the Stripe API; the agent answers five customer tickets with known correct answers.
\item \textbf{C5, CI triage}$^\dagger$: six tools of a simulated MCP server following the GitHub REST API
  description; the agent explains the week's CI failures on \texttt{main} and links them to issues and pull
  requests, again with known correct answers.
\end{itemize}
$^\dagger$Simulated payloads are generated from the OpenAPI schemas with a seeded generator; because we design
them, C4 and C5 are reported separately and never pooled with the real-API cases. Inputs name only the author's
own domain; a gate in the tool wrapper rejects any other domain for page parsing, ranking targets and
\texttt{site:} queries.

\paragraph{Study 1.} \emph{Phase A}: eight live runs per case, used for usage, leave-one-out projections,
learning source and transfer between C1 and C2 (on the tools they share). \emph{Phase B}: the third run of each
case, chosen before looking at the data, is replayed up to the final answer, changing only the tool messages:
raw, prudent, aggressive, and prudent plus structural compression; two runs per arm; projections learned from the
other seven runs. Headroom (\texttt{headroom-ai} \HeadroomVersion) is applied to the same conversations as an illustrative comparison of
tokens only, in the configuration chosen below, without any evaluation of the answers it would lead to. It ran
\emph{stateless}: each conversation compressed once, so its per-tool learning had no history of these tools, while
the projection learned from seven runs. It was configured to compress everything, with no protection of recent
messages and a target ratio of 0.2 (an aggressive setting, not one tuned for these tasks); its input was the
conversation's user message, tool calls and tool messages, without the system prompt, with the model parameter set to
\texttt{gpt-4o}. \emph{Phase C}: live
runs on C1 and C4 with the prudent projection and \texttt{expand}, against raw runs, three per arm.

\paragraph{Study 2.} The Phase B conversations are replayed again with the safeguards on (arm \texttt{v02}, and
\texttt{v02} plus structural compression) against the raw arm, five runs per arm.

\paragraph{Measures.} \emph{Gates} are binary: the answer parses, has the required structure, and the facts it
attributes to tools match the tool outputs stored in the trace. As implemented they are narrower than the
protocol's wording: C1 checks each reported volume against the volumes the run's outputs give for that keyword,
and skips a keyword that has none; C2 checks that each cited URL was returned by the search tool, not its title;
C3 checks each work's title through its identifier, and that a declared identifier exists; C4 checks that every
object identifier cited exists; C5 checks run, job and issue numbers. Appendix~\ref{app:supplementary} measures
what this leaves unchecked. In Phases A and B and in Study 2 the stored outputs are the raw ones, whatever the agent
was shown; in Phase C, where the live loop stores what the agent received, they are the projected outputs and what
\texttt{expand} returned (Section~\ref{sec:phasec}). \emph{Thermometers} are graded: editorial metrics for C1,
coverage metrics for C2 and C3, and, for C4 and C5, keyword tests for the diagnoses and links whose correct answers
are known. Results are reported as medians with ranges (Appendix~\ref{app:supplementary}); with two to five runs
per arm no significance test is claimed. The pre-registered rule for quality (H2) is relative: an arm \emph{meets}
it when it passes every gate in at least as many runs as the raw arm and its median on every thermometer lies
within the raw arm's range. It says nothing about absolute quality: where no raw run passes the gates, an arm meets
it with none passing either.
